\documentclass[11pt]{article}
\usepackage[margin=2.5cm]{geometry}
\usepackage[utf8]{inputenc}
\usepackage[T1]{fontenc}
\usepackage{enumitem}
\usepackage{titlesec}
\usepackage{xcolor}
\usepackage{hyperref}
\usepackage{parskip}

\definecolor{headcol}{RGB}{30,60,110}
\titleformat{\section}{\Large\bfseries\color{headcol}}{}{0pt}{}
\titleformat{\subsection}{\large\bfseries\color{headcol}}{}{0pt}{}
\titlespacing*{\section}{0pt}{1.4em}{0.7em}
\titlespacing*{\subsection}{0pt}{1.2em}{0.5em}

\newcommand{\code}[1]{\texttt{#1}}

\hypersetup{colorlinks=true, linkcolor=headcol, urlcolor=headcol}

\title{\vspace{-1.5cm}\textbf{Thesis Project Summary}\\[4pt]
\large Unsupervised Clustering of Alzheimer's/Dementia Patients from Brain MRI}
\author{}
\date{}

\begin{document}
\maketitle
\vspace{-2.2cm}

\section*{Goal}
Using the OASIS cross-sectional dataset (416 patients, 5 MRI scan types per patient), the
question is whether \textbf{image-derived features alone, or combined with clinical variables,
produce patient clusters that align with clinical dementia severity (CDR)} --- without ever using
CDR as a clustering input. CDR (0 = healthy, 0.5 = very mild, 1 = mild, 2 = moderate) is known for
only 235 of the 416 patients and is used strictly for \textbf{post-hoc evaluation} of whatever
clusters come out.

\section*{Feature extraction}
Each patient has five MRI planes (coronal, atlas-registered sagittal and transverse, a
brain-masked transverse, and a native/unregistered sagittal). Each image is padded, resized, and
converted to HOG features (9 orientations, 8$\times$8 cells, 2$\times$2 blocks $\rightarrow$ 8{,}100
features/image, 40{,}500 total per patient). Per image type, PCA is applied with a two-pass
strategy that finds the minimum component count needed for $\sim$59\% variance in the hardest
type, then refits every type to that same number of components (110), so all five image types are
directly comparable dimensionally.

\section*{Clustering framework}
Clustering uses a family of robust, mixed-type distance metrics and a fast k-medoids engine
(\code{robust\_mixed\_dist} / \code{db\_robust\_clust}, developed by our advisor Fabio
Scielzo-Ortiz). This is the same toolkit used in his Reddit/Gaza-conflict polarization paper ---
there it clusters LLM-derived discursive variables extracted from text; here it clusters
PCA-reduced HOG features extracted from MRI images. The core methods:
\begin{itemize}[leftmargin=1.4em]
  \item \textbf{Robust Mahalanobis} distance for single quantitative blocks
  \item \textbf{Generalised Gower} distance for mixing quantitative + binary + categorical clinical variables
  \item \textbf{Related Metric Scaling (RelMS)} for combining several \emph{typed blocks} (e.g.\ five
  separate image types), each rescaled by its own geometric variability before being combined, so
  no single block dominates purely because it has more columns
  \item \textbf{SampleDistClustering / FasterPAM k-medoids}, which subsamples for tractability and
  returns real patients as cluster medoids (not synthetic centroids)
\end{itemize}

\section*{What was actually run}
\begin{enumerate}[leftmargin=1.4em, itemsep=0.4em]
  \item \textbf{Per-image-type clustering} (Euclidean and Robust Mahalanobis, $k=3/4$, each of the
  5 image types alone) --- best single-type result was \code{masked\_tra} (Euclidean, $k=3$):
  ARI~$=0.060$, NMI~$=0.085$, the strongest CDR alignment found anywhere in the project.

  \item \textbf{Merged 5-image clustering} (all 550 PCA features concatenated, Euclidean vs.\
  Robust Mahalanobis, $k=3/4$) --- ARI stayed in $[-0.007, 0.013]$, NMI in $[0.02, 0.05]$.

  \item \textbf{Clinical + imaging via Generalised Gower} (all-images-merged, each single image
  type, and clinical-only, all combined with 8 clinical variables) --- revealed an anomaly: every
  one of these variants (all-images+clinical, cor+clinical, gfc\_sag+clinical, \ldots, and even
  the clinical-only baseline) produced numerically \textbf{identical} results (silhouette
  $\approx 0.328$, ARI $=0.016$, NMI $=0.022$ at $k=3$), down to the same cluster sizes. This
  strongly suggests the clinical block was dominating the Generalised Gower distance so completely
  that the 110--550 image PCA columns contributed essentially zero signal.

  \item \textbf{Two-image concatenation} (masked\_tra + cor, Euclidean) --- ARI still near zero.

  \item \textbf{RelMS distance} (Related Metric Scaling), introduced specifically to fix the
  block-dominance problem from step 3, applied to: all 5 images merged; all 5 images + clinical
  (complete cases, $n=216$); each single image + clinical; and the two most promising two-image
  combinations (\code{masked\_tra+cor}, \code{masked\_tra+gfc\_tra}), with and without clinical
  data. Best result: \code{masked\_tra+cor+clinical}, $k=4 \rightarrow$ ARI~$=0.041$,
  NMI~$=0.041$ --- still weak.

  \item \textbf{Image-type redundancy analysis} --- correlated the pairwise-patient-distance
  matrices of the five image types against each other. Most redundant pair: \code{cor}
  $\leftrightarrow$ \code{gfc\_sag} ($r=0.60$). Most complementary/independent pair:
  \code{masked\_tra} $\leftrightarrow$ \code{sbj\_sag} ($r=0.016$) --- meaning the native
  (unregistered) sagittal images carry the most unique structural signal, distinct from
  everything else.

  \item \textbf{Leave-one-image-out sweep} --- systematically removed each image type from the
  RelMS 5-image blend to see which types help vs.\ hurt CDR alignment. Removing \code{masked\_tra}
  or \code{sbj\_sag} tended to \emph{improve} ARI/NMI slightly (they add noise more than signal in
  the 5-way blend), while removing \code{gfc\_sag} or \code{gfc\_tra} tended to hurt it --- but all
  deltas were small ($\leq 0.01$), i.e.\ within noise.

  \item \textbf{Full cluster profiling} (5 images + clinical, complete cases $n=216$) --- beyond
  ARI/NMI, this notebook profiled each cluster's mean age, MMSE, brain volume (eTIV/nWBV),
  education, SES, and CDR/sex breakdown, with visualizations (boxplots, stacked CDR bars, a
  z-scored cluster-mean heatmap, 2D/3D PCA scatter, and effect-size ($\eta^2$/Cram\'er's $V$)
  rankings of which clinical features actually separate the clusters). It also surfaced a second
  anomaly: at $k=3/4$ on a related Tucker configuration, Robust Mahalanobis produced degenerate
  singleton-outlier clusters (e.g.\ sizes $\{1, 414, 1\}$) with a deceptively high silhouette
  ($\sim 0.78$) --- an artifact of isolating 1--2 outliers, not real structure.

  \item \textbf{Per-block distance sensitivity (newest addition)} --- following an update from the
  advisor (\code{robust-mixed-dist} 0.1.27) allowing a \emph{different} distance per quantitative
  block within RelMS, this was tested on the same 5-image + clinical setup: keeping the five image
  blocks on Euclidean while switching only the clinical block to Robust Mahalanobis, switching
  \emph{all} quantitative blocks to Robust Mahalanobis, and switching only \code{masked\_tra}. All
  four configurations were also compared against four alternative distance \emph{functions}
  (Minkowski, Canberra, Pearson, plain Mahalanobis). Result: the distance choice does meaningfully
  change \emph{who ends up in which cluster} (pairwise ARI between distance variants ranged from
  0.11 to 0.98), but CDR alignment never moves outside the noise band already seen everywhere else
  (best ARI vs.\ CDR $\approx 0.03$, best NMI $\approx 0.02$) --- so the distance metric was not the
  bottleneck.
\end{enumerate}

\section*{Central finding}
Across every method tried --- single image type, merged images, Generalised Gower, RelMS,
two-image combinations, leave-one-out sweeps, and multiple distance-function variants ---
\textbf{clustering does not meaningfully recover CDR groups}. ARI stays roughly in
$[-0.02, +0.06]$ and NMI in $[0.01, 0.10]$ throughout, i.e.\ barely above what random cluster
assignment would produce. The best-performing single result across the whole project was
\code{masked\_tra} alone with Euclidean distance at $k=3$ (ARI $=0.060$). This is a genuine
negative result, not a bug: the same features that separate patients geometrically (image
texture/shape via HOG+PCA) simply don't line up with the axis clinicians use (CDR severity), and
swapping distance metrics, block weightings, or image combinations doesn't change that.

\section*{Two things worth flagging as anomalies (not fully resolved)}
\begin{enumerate}[leftmargin=1.4em, itemsep=0.3em]
  \item In the Generalised Gower experiments (step 3), all image+clinical variants were
  numerically identical to the clinical-only baseline --- pointing to a likely scaling issue where
  the clinical block dominates that particular distance formulation.
  \item In one Tucker-decomposition configuration, Robust Mahalanobis produced degenerate
  singleton clusters with an artificially high silhouette score --- a reminder that silhouette
  alone can be misleading when a method isolates a couple of extreme outliers rather than finding
  real structure.
\end{enumerate}

\end{document}
